%% AMS-LaTeX Created with the Wolfram Language : www.wolfram.com

\documentclass[12pt]{article}
\usepackage[letterpaper,margin=1in]{geometry}
\usepackage{amsmath, amssymb, graphics, setspace}
\usepackage{bm}
\usepackage{graphicx} \graphicspath{ {./figs/} }
\usepackage{tcolorbox}
\usepackage[framed,numbered,autolinebreaks,useliterate]{mcode}
\newcommand{\mathsym}[1]{{}}
\newcommand{\unicode}[1]{{}}

\usepackage[english]{babel}
\newtheorem{theorem}{Theorem}

\usepackage{caption}
\usepackage{subcaption}

\usepackage{tabulary}
\usepackage{multirow}
\usepackage{lscape}

\usepackage{hyperref}
\hypersetup{
    colorlinks=true,
    linkcolor=blue,
    filecolor=magenta,      
    urlcolor=cyan,
    pdftitle={Overleaf Example},
    pdfpagemode=FullScreen,
    }

\def\mathbi#1{\textbf{\em #1}}
\newcommand{\field}[1]{\mathbb{#1}}
\newcommand{\Rset}{\field{R}}
\newcommand{\Zset}{\field{Z}}
\newcommand{\Qset}{\field{Q}}
\newcommand{\Nset}{\field{N}}

\begin{document}
\doublespacing	

\title{CE 401/5501: Applied AI \\ \large Topic 14: Machine Learning for Regression}
\author{Dr. ZhiQiang Chen}
\date{}
\maketitle

\section{Review of Regression}

\subsection{What Is Regression?}
Regression analysis studies the mapping
\[
  \bm{x} \;\longmapsto\; y
\]
where \(\bm{x}\in\Rset^p\) is a vector of one or more input (predictor) variables and \(y\in\Rset\) is a continuous output.  In engineering and the sciences, regression underlies:
\begin{itemize}
  \item \textbf{Prediction:} forecasting future behavior (e.g.\ loads, flows, costs).
  \item \textbf{Modeling:} capturing system relationships (e.g.\ stress–strain).
  \item \textbf{System Identification:} inferring process parameters from measurements.
\end{itemize}
By contrast, \emph{classification} assigns inputs to a finite set of categories.


\noindent
\colorbox{green!10}{%
  \parbox{\textwidth}{%
    \textbf{Practical Tip:} Always begin with a simple linear regression model to establish a baseline. Exploring and transforming your inputs (e.g., via log or power transforms) often yields substantial improvements before moving to complex methods.
  }%
}

\subsection{Linear Least Squares Regression}
Suppose we posit a linear model
\[
  y \;=\; f(\bm{x};\bm{w})
  \;=\;\bm{w}^T\bm{x}
  \quad\text{or with intercept}\quad
  y = w_0 + \bm{w}^T\bm{x}.
\]
Given \(N\) observations \(\{(\bm{x}_i,y_i)\}_{i=1}^N\), collect
\(\bm{y}=(y_1,\dots,y_N)^T\) and \(X=[\bm{x}_1\,\dots\,\bm{x}_N]^T\).  The ordinary least‐squares criterion is
\[
  J(\bm{w})
  \;=\;
  \bigl\|\bm{y} - X\,\bm{w}\bigr\|^2
  \;=\;
  \sum_{i=1}^N \bigl(y_i - \bm{w}^T\bm{x}_i\bigr)^2.
\]
Because \(J(\bm{w})\) is a convex quadratic function, it has the closed‐form minimizer
\[
  \hat{\bm{w}}
  =
  (X^T X)^{-1} X^T \bm{y},
\]
provided \(X^T X\) is invertible.

\subsection{Nonlinear Regression and Basis Expansions}
Not all relationships are linear in the raw predictors.  Two common remedies:
\begin{enumerate}
  \item \emph{Transform the response or inputs} so that a linear relationship emerges.
  \item \emph{Basis‐function expansion:} model
    \(\;y = \sum_{j=1}^M w_j\,\phi_j(\bm{x})\), where \(\{\phi_j\}\) are nonlinear basis functions.
\end{enumerate}

\paragraph{Example 1: Power‐Law Regression.}
\[
  y = a\,\prod_{j=1}^p x_j^{b_j}
  \quad\Longrightarrow\quad
  \ln y = \ln a \;+\;\sum_{j=1}^p b_j\,\ln x_j,
\]
so a log–log plot followed by linear regression estimates \(\ln a\) and each \(b_j\).

\paragraph{Example 2: Rectangular Hyperbola.}
\[
  y
  = 
  \frac{\beta_1\,x}{\beta_2 + x},
\]
which is intrinsically nonlinear but can be fitted by nonlinear least squares or linearized via the reciprocal transform.

\subsection{Engineering Examples}

\paragraph{Cost Prediction in Construction.}  
Model total project cost \(y\) as
\[
  y 
  = 
  \beta_0 
  + \beta_1\,x_{\text{mat}}
  + \beta_2\,x_{\text{lab}}
  + \beta_3\,x_{\text{dur}}
  + \beta_4\,x_{\text{size}},
\]
where
\(\;x_{\text{mat}}\)=material cost,\;
\(x_{\text{lab}}\)=labor cost,\;
\(x_{\text{dur}}\)=duration (weeks),\;
\(x_{\text{size}}\)=area (sq.\ ft).  Fit by minimizing
\[
  J(\boldsymbol{\beta})
  = 
  \sum_{i=1}^N 
    \bigl(
      y_i 
      - 
      \beta_0 
      - 
      \sum_{j=1}^4 \beta_j\,x_{ij}
    \bigr)^2.
\]

\paragraph{Concrete Modulus Prediction (ACI 318).}  
Empirical formula for Young’s modulus \(E_c\):
\[
  E_c 
  = 
  33\,w_c^{1.5}\,\sqrt{f'_c},
  \quad
  f'_c=\text{compressive strength (MPa)},\quad
  w_c=\text{unit weight (lb/ft}^3).
\]
Log‐transform yields
\[
  \log E_c 
  = 
  \alpha_0 
  + 
  \alpha_1\,\log w_c 
  + 
  \alpha_2\,\log f'_c,
\]
reducing to a straight‐line fit in \((\log w_c,\log f'_c)\) space.

\section{Logistic Regression}

\subsection{Logistic Regression as Classification}
Logistic regression is a \emph{binary} classification method that maps an input vector \(\bm{x}\in\Rset^p\) to an estimated probability that the class label \(y\in\{0,1\}\) equals 1:
\[
  \hat p(\bm{x}) \;=\; P(y=1 \mid \bm{x})
  \;=\;
  \sigma\bigl(\bm{w}^T\bm{x} + b\bigr),
\]
where
\[
  \sigma(u) \;=\;\frac{1}{1+e^{-u}}
\]
is the logistic (sigmoid) function.

\begin{figure}[h!]
  \centering
  \includegraphics[width=0.6\textwidth]{figs/sigmoid.pdf}
  \caption{The sigmoid activation function \(\sigma(u)=1/(1+e^{-u})\).}
  \label{fig:sigmoid}
\end{figure}

As shown in Figure~\ref{fig:sigmoid}, \(\sigma(u)\) smoothly maps \(\Rset\) to \((0,1)\).  The decision rule classifies \(\bm{x}\) as
\[
  \hat y = 
  \begin{cases}
    1, & \hat p(\bm{x}) \ge 0.5,\\
    0, & \text{otherwise},
  \end{cases}
\]
which corresponds to the linear decision boundary \(\bm{w}^T\bm{x} + b = 0\).

\subsection{Parameter Estimation via Maximum Likelihood}
Given training data \(\{(\bm{x}_i,y_i)\}_{i=1}^N\), we choose \((\bm{w},b)\) to maximize the likelihood
\[
  L(\bm{w},b)
  = 
  \prod_{i=1}^N
    \hat p(\bm{x}_i)^{y_i}
    \bigl(1-\hat p(\bm{x}_i)\bigr)^{1-y_i}.
\]
Equivalently, we minimize the negative log‐likelihood (the \emph{cross‐entropy} loss):
\[
  J(\bm{w},b)
  = 
  -\sum_{i=1}^N
    \bigl[
      y_i\ln\hat p(\bm{x}_i)
      + (1-y_i)\ln\bigl(1-\hat p(\bm{x}_i)\bigr)
    \bigr].
\]
Since \(J(\bm{w},b)\) is convex in \((\bm{w},b)\), standard optimization methods (e.g.\ gradient descent, Newton’s method) find the global minimum.

\paragraph{Gradient of the Loss}
\[
  \frac{\partial J}{\partial \bm{w}}
  = 
  \sum_{i=1}^N \bigl(\hat p(\bm{x}_i) - y_i\bigr)\,\bm{x}_i,
  \quad
  \frac{\partial J}{\partial b}
  = 
  \sum_{i=1}^N \bigl(\hat p(\bm{x}_i) - y_i\bigr).
\]

\subsection{Interpretation of the Logit}
Define the \emph{log‐odds} (logit)
\[
  \log\frac{\hat p(\bm{x})}{1-\hat p(\bm{x})}
  = \bm{w}^T\bm{x} + b.
\]
Thus logistic regression is a \emph{linear} model on the log‐odds scale.

\subsection{Regularization}
To prevent overfitting, one often adds a penalty term:
\[
  J_{\lambda}(\bm{w},b)
  = 
  J(\bm{w},b)
  + \frac{\lambda}{2}\|\bm{w}\|^2
  \quad\text{(L2 regularization)},
\]
which keeps the optimization convex and shrinks weights toward zero.

\section{Machine Learning for Regression}

\subsection{Multi‐Layer Perceptron (MLP)}

A multi‐layer perceptron is a feedforward neural network composed of one or more hidden layers.  For regression, the final layer is linear and produces a continuous output:
\[
  \hat y(\bm{x})
  \;=\;
  \bm{w}_L^T\,\bm{h}_{L-1}(\bm{x}) \;+\; b_L,
\]
where \(\bm{h}_{L-1}\) is the activation of the last hidden layer.  Training minimizes the mean‐squared error (MSE):
\[
  J(\Theta)
  = 
  \frac{1}{N}\sum_{i=1}^N \bigl(y_i - \hat y(\bm{x}_i)\bigr)^2.
\]

\subsubsection{Example: MLP Regression}

\begin{figure}[h!]
  \centering
  \includegraphics[width=0.6\textwidth]{figs/sin-mlp.pdf}
  \caption{MLP fit to the nonlinear function \(y = \sin(x)\).}
  \label{fig:sin_mlp}
\end{figure}

\noindent
\colorbox{blue!10}{%
  \parbox{\textwidth}{%
    \textbf{Universal Approximation Theorem.}  
    A single‐hidden‐layer MLP with sufficiently many neurons can approximate any continuous function on a compact domain to arbitrary precision:
    \[
      \Bigl|\,f(\bm{x}) \;-\; \sum_{i=1}^M \alpha_i\,\sigma(\bm{w}_i^T\bm{x} + \theta_i)\Bigr|
      < \varepsilon
      \quad\forall\,\bm{x}\in[0,1]^n.
    \]
    Here \(\sigma\) is an activation function such as sigmoid, tanh, or ReLU.
  }%
}

\vspace{1em}
\noindent
\colorbox{yellow!20}{%
  \parbox{\textwidth}{%
    \textbf{Practical Challenges.}\  
    Although universal approximation holds in theory, in practice one must contend with:
    \begin{itemize}
      \item Optimization difficulties (local minima, saddle points).  
      \item Overfitting without sufficient data or regularization.  
      \item Choosing appropriate architecture (number of layers/neurons).
    \end{itemize}
  }%
}

\subsection{Support Vector Regression (SVR)}

Support vector regression uses the kernel trick together with an \(\varepsilon\)-insensitive loss that ignores small errors:
\[
  L_{\varepsilon}(y,f(\bm{x})) 
  = 
  \max\{0,\,|y - f(\bm{x})| - \varepsilon\}.
\]
The resulting predictor is
\[
  f(\bm{x})
  = 
  \sum_{i\in\mathrm{SV}} (\alpha_i - \alpha^*_i)\,K(\bm{x}_i,\bm{x}) + b,
\]
where only support vectors (\(\mathrm{SV}\)) with nonzero \(\alpha\) terms contribute.

\begin{figure}[h!]
  \centering
  \includegraphics[width=0.6\textwidth]{figs/cos-svr.pdf}
  \caption{SVR applied to noisy samples of \(y=\cos(x)\).}
  \label{fig:cos_svr}
\end{figure}

\subsection{Decision Tree Regression}

Decision trees partition the input space into axis‐aligned regions \(R_m\).  For regression, splits are chosen to minimize within‐node variance:
\[
  \min_{j,t}
  \Biggl\{
    \sum_{i:x_i\in R_{\text{left}}} (y_i - \bar y_{\text{left}})^2
  +
    \sum_{i:x_i\in R_{\text{right}}} (y_i - \bar y_{\text{right}})^2
  \Biggr\}.
\]
Each leaf node \(R_m\) predicts the average of its responses:
\[
  \hat y(R_m) = \frac{1}{|R_m|}\sum_{i: x_i\in R_m} y_i.
\]

\begin{figure}[h!]
  \centering
  \includegraphics[width=0.45\textwidth]{figs/dt_visualization.png}
  \caption{Partitioning of the domain for \(y=x^2\).}
  \label{fig:dt_partition}
\end{figure}
\begin{figure}[h!]
  \centering
  \includegraphics[width=0.45\textwidth]{figs/dt-fitting.pdf}
  \caption{Resulting piecewise‐constant fit to noisy \(y=x^2\) data.}
  \label{fig:dt_fit}
\end{figure}

\noindent
\colorbox{orange!10}{%
  \parbox{\textwidth}{%
    \textbf{Hyperparameter Tuning:} Use cross‐validation to select key settings—network size and regularization strength in MLPs; \(\varepsilon\), \(C\) (margin penalty) in SVR; and tree depth or pruning criteria in decision trees.
  }%
}

\section{Summary}

Regression techniques map inputs to continuous outputs. Start with linear least squares as an interpretable, convex baseline. For nonlinear patterns, choose among:
\begin{itemize}
  \item \textbf{MLP:} Flexible universal approximator, requires careful regularization and architecture selection.
  \item \textbf{SVR:} Sparse kernel methods with \(\varepsilon\)-insensitive loss, balancing margin width (\(\varepsilon\)) and penalty (\(C\)).
  \item \textbf{Decision Trees:} Interpretable piecewise‐constant fits, controlled via depth and pruning to avoid overfitting.
\end{itemize}

\noindent
Balancing bias vs.\ variance—and interpretability vs.\ complexity—is key. Always validate with held‐out data and communicate model choices clearly for engineering decision‐making.






Ask ChatGPT

\end{document}